% A unit baseline and a direct guard preserve the fixed-four alphabet.
\section{A coefficient baseline of one}\label{sec:unit-center}

The coefficient baseline in Theorem~\ref{thm:best99} can be reduced
from two to one. Copies remove the negative square coefficient in
the unit test, while a direct guard excludes a zero homogenizing
coordinate before any copy relation is used. The resulting
coefficient is
\begin{equation}\label{eq:unit-center-block}
 \Omega=\lambda-e>0,\qquad
 S_3=\theta\lambda q-\Omega\mathcal C^2=\sigma>0,
 \qquad\mathcal C=x+g,\quad\theta=B-4.
\end{equation}
The offset, geometric equation, packed upper bound and all Pell
equations are retained. Computing $3\lambda$ directly now saves
one addition, giving 98 operations.

\subsection{Copied unit tests and a direct guard}

Keep the paired primitive rows of Section~\ref{sec:half-center}:
each multiplication, addition, copy, final equality and zero
residual $F$ occurs together with $-F$, with targets
$2F+\delta^2$. Operands are distinct copied coordinates, outputs
are fresh, and circuit coordinates are distinct from $\delta$.

For each unit gate $X=1$, introduce a fresh nonnegative coordinate
$X'$ and the paired copy rows $(X-X')\delta$ and $(X'-X)\delta$.
Replace its two unpaired unit rows by
\begin{equation}\label{eq:unit-center-unit}
 F_{X,1}=X\delta-\delta^2,\qquad
 F_{X,2}=X\delta-XX'.
\end{equation}
All four rows still use targets $2F+\delta^2$. The negative product
$XX'$ has distinct factors, so its divided target coefficient is
$-1$, in place of the $-2$ of the old square $X^2$.

Give the fresh guard coordinate $u$ the direct target
\begin{equation}\label{eq:unit-center-guard}
 G_{\mathrm g}=2u\delta-x^2+\delta^2.
\end{equation}
This target is not required to be $2F+\delta^2$ for an integer
polynomial $F$. It will exclude $\delta=0$ directly. Replacing
$x^2$ by a copied product would not do this: its homogenized copy
equation would impose no condition when $\delta=0$.
Keep the special target $G_0=\delta^2$ first.

Every target coefficient, divided by its monomial coefficient in
$\mathcal C(T)^2$, now belongs to
\begin{equation}\label{eq:unit-center-alphabet}
 \{-1,0,1\}.
\end{equation}
Paired primitive rows have cross-term quotients $\pm1$ and added
$\delta^2$ coefficient $1$. The first unit target has coefficients
$1,-1$ after division; the second has cross-term quotients $1,-1$
and square coefficient $1$. The direct guard has cross-term quotient
$1$, coefficient $-1$ on $x^2$, and coefficient $1$ on $\delta^2$.
Thus all these divisions are integral.

\subsection{The explicit admissible index}

Apply the support construction to the entire new finite list,
including all unit-copy coordinates and paired copy rows. Let $m$
be the number of positive-weight coordinates and $s$ the number of
targets. The input $x$ has weight zero and $\delta$ has weight one.
Set
\begin{align*}
 v_i&=3^i\quad(0\le i<m),& M&=3^{m-1},& d_0&=4M+1,\\
 t_j&=(s+1)d_0+2M+jd_0\quad(0\le j<s),\\
 t_*&=2sd_0+2M,& K_0&=t_*+1,&c_*&=m+s+1.
\end{align*}
The special target occupies $t_0$. Include a dummy coordinate at
every target position and choose a power of two $L>3K_0+2$.
The complementary coefficient polynomial $\mathcal D$ has
coefficients in \eqref{eq:unit-center-alphabet} and satisfies
\[
 [T^{t_j}]\,\mathcal D(T)\mathcal C(T)^2=G_j.
\]
The base-three monomial uniqueness and disjoint row intervals prove
this identity exactly; $2t_0-2M>t_*$ excludes every dummy contribution.
The support of $\mathcal D$ lies above every low coordinate and
below $K_0$. Through $t_0$ its only term is $+T^{t_0-2}$, since
every later row starts above $t_0$.

Define
\begin{align}
 \ell_0(T)&=\sum_{i=0}^{m-1}T^{v_i}+\sum_{j=0}^{s-1}T^{t_j},\notag\\
 e_0(T)&=\sum_{j=0}^{K_0-1}T^j+\mathcal D(T),&
 V&=\ell_0(4)+e_0(4)4^L.\label{eq:unit-center-index}
\end{align}
The digits of $e_0$ are in $\{0,1,2\}$ and its unit digit is one.
In particular, $e_0>0$, while the packed digit polynomial has all
digits below four and degree below $2L$. Choose a power of two
\begin{equation}\label{eq:unit-center-H}
 H>\max\{2\,4^{2L+1},\ 4^{t_*+3}\,4c_*^2,\ 3L,\ 16\}.
\end{equation}
Finally supply $T_L=\psi_4(L)$. The admissible index is
$(V,H,T_L)$; $L$ is the underlying exponent defining it, not an
additional system input. Both $V$ and $H$ are rebuilt from the new
layout and baseline. All choices depend only on the represented
system, independently of the queried positive input $x$.

\subsection{Establish the Pell hypotheses before decoding}

Besides \eqref{eq:unit-center-block}, retain
\begin{align*}
 B&=Hb^2,& b&=x+\beta,&n&=q^8,\\
 \lambda(B-1)&=q^2-1,&\theta+4&=B,\\
 S_2&=\ell+eq=V+t\theta,&S_2+\alpha&=q^2,\\
 S&=g+q^2(S_2+q^2S_3),\\
 T+1&=q^2(1+\theta\lambda)-(b-1)\ell+\theta\ell q^4,\\
 r&=S(n^2-n)+(T+1)(n^2-1).
\end{align*}
Every supplied unknown remains positive. The packed bound gives
$e<q$ and $\ell<q^2$, and geometry gives $B\le q^2$.
The bound on $H$ gives $q>4b\ge8$ and $3L\le B\le n$.
Since $\Omega\ge1$, positivity of the new block gives directly
\[
 \mathcal C^2<\theta\lambda q<B\lambda q<3q^3<q^4.
\]
Thus $g<\mathcal C<q^2$ and $0<S_3<q^4$. The first-block borrow
is fewer than $b$, while $\theta\lambda>b$, so the unchanged
width-$(2,2,4)$ estimates give
\begin{equation}\label{eq:unit-center-ranges}
 0<S<n,\quad0<T+1\le n,\quad
 n\le r\le2n^3-2n^2<2n^3.
\end{equation}
These are direct bounds for the new coefficient; no comparison
with an old value of $S_3$ is assumed.

Write $U=wn^2$, $Y_P=sn^2$, $a=Y_P(U+1)$, $A=a+4$ and $J=2r+1$.
Then $n,r\ge64$, $U,Y_P\ge n^2$, $UY_P\ge n^4>r+1$, and
$a>n^4>J$. The first index congruence and positive interval give
$c>J$. Thus $A>1$, $1<J<c$ and odd $J$ are available, as are
$L\ge2$ and $3L\le B\le n\le r$.

All Pell source equations of Theorem~\ref{thm:best99} are unchanged.
Their proof therefore applies in its original order. First the
doubled-index auxiliary norm identifies $c=\psi_A(J)$,
$d=\chi_A(J)$. The first norm and index give
$k=\psi_{2UY_P^2+1}(r+1)$. The lower ratio estimate gives
$Y_P\ge U^r$ and $a>U^{r+1}$, before the first exponent equation
yields $U=4^J$. In particular $a>4^{3J}$.
The positive gap and the equation $\kappa=T_L+\Delta a$ now
identify the second Pell index as $L$: its two residues
$\psi_4(t)$ and $\psi_4(L)$ lie strictly between zero and $a$.
The second exponent comparison uses
\[
 B^{3L}\le B^B\le n^n\le U^r<A,
 \qquad q^3\le n^n\le U^r<A,
\]
and yields $q=B^L$. The upper ratio and binomial rounding arguments
then give
\begin{equation}\label{eq:unit-center-decoded-powers}
 b,B,q\text{ are powers of two},\qquad
 q=B^L,\qquad n^2\mid\binom{2r}{r}.
\end{equation}
Every internal coefficient test still awaits decoding.

The weaker coefficient requires a revised post-Pell estimate. Since
$L\ge2$ and $B>16$,
\[
 \lambda\ge B^{2L-1}\ge Bq>2q,
 \qquad q=B^L\ge B^2>2B.
\]
Therefore $e<q$ implies $\Omega=\lambda-e>\lambda/2$, and
\eqref{eq:unit-center-block} yields
\begin{equation}\label{eq:unit-center-C-bound}
 \mathcal C^2<2\theta q<2Bq<q^2,
 \qquad g<\mathcal C<q.
\end{equation}

\subsection{Canonical codes and the smaller coefficient plateau}

The first two masks, congruence and bound are unchanged. Their
borrow-removal proof gives $d=\lfloor(b-1)\ell/q^2\rfloor=0$:
the middle mask first implies
$\ell\le S_2<4q^2/(B-1)+b$, which forces $(b-1)\ell<q^2$.
The base-four evaluation lemma applies to
\eqref{eq:unit-center-index}; it permits digits in $\{0,1,2\}$
and does not require the digit three. Consequently
\[
 S_2=\ell_0(B)+e_0(B)q,\qquad
 \ell=\ell_0(B)+a_{\mathrm{al}}q,\qquad
 e=e_0(B)-a_{\mathrm{al}},\qquad
 0\le a_{\mathrm{al}}<e_0(B).
\]
Using \eqref{eq:unit-center-C-bound}, the first mask decodes every
coordinate below $b$, with input $x$ at the unit position and dummy
digits only at row positions. Thus
\[
 e,\mathcal C<B^{K_0},\qquad e\mathcal C^2<q,
 \qquad A_{\mathcal C}=\mathcal C(1)^2<(c_*b)^2.
\]
Since $x,g>0$, one has $\mathcal C(1)\ge2$ and
$A_{\mathcal C}\ge4$ before recovering $\delta$.

Now the polynomial $\lambda\mathcal C^2$ has coefficients at most
$J_{\mathcal C}=A_{\mathcal C}<B/4$, and its digits from $L$
through $L+K_0$ all equal $J_{\mathcal C}$. Write
\[
 \lambda\mathcal C^2=qV_1+V_0,\qquad0\le V_0<q.
\]
Exact division gives
\[
 \left\lfloor\frac{S_3}{q}\right\rfloor
 =\theta\lambda-V_1+\epsilon,\qquad
 \epsilon=\left\lfloor
 \frac{e\mathcal C^2-V_0}{q}\right\rfloor\in\{-1,0\}.
\]
There is no borrow, since the digits of $\theta\lambda$ are $B-4$
throughout this interval. The first digit of $S_3$ in the window is
$B-4-J_{\mathcal C}+\epsilon$, and every subsequent one is
$B-4-J_{\mathcal C}$. Binary nonoverlap therefore bounds every
digit of $X=(B-4)a_{\mathrm{al}}$ by
\[
 J_{\mathcal C}+4=A_{\mathcal C}+4\le4A_{\mathcal C}.
\]
For its digit polynomial $F_X$,
\[
 B-4\mid F_X(4),\qquad
 0\le F_X(4)\le
 4A_{\mathcal C}\frac{4^{K_0+1}-1}{3}<\frac B{12}<B-4.
\]
The unchanged bound \eqref{eq:unit-center-H} supplies the strict
inequality. Hence $X=0$ and $a_{\mathrm{al}}=0$, so both codes
are canonical before any circuit row is tested.

\subsection{Recover the special coordinate before the copies}

Below $K_0$, the raw coefficients of $S_3$ are now those of
$\mathcal D(T)\mathcal C(T)^2$. Each has absolute value at most
$A_{\mathcal C}<B/4$, so incoming carries are zero or minus one.
At a target raw value $v$, the mask $B-4$ is equivalent to
\begin{equation}\label{eq:unit-center-target-test}
 0\le v+\epsilon\le3,\qquad\epsilon\in\{-1,0\}.
\end{equation}
The first special target has incoming carry zero: its single
positive coefficient precedes all later row intervals, making
every raw coefficient through that target nonnegative and below
$B$. Thus $0\le\delta^2\le3$, so $\delta\in\{0,1\}$.

At $\delta=0$, the direct guard \eqref{eq:unit-center-guard} has
raw value $-x^2<0$, which fails
\eqref{eq:unit-center-target-test}. Hence $\delta=1$ without any
copy assumption. Ordinary targets now test $F\in\{0,1\}$, and
paired rows force $F=0$. In particular, the paired unit copies give
$X'=X$. The two rows \eqref{eq:unit-center-unit} then impose
\[
 X-1\in\{0,1\},\qquad X-X^2\in\{0,1\},
\]
forcing $X=1$. Every original circuit gate and equality is recovered
at the actual input. The guard is auxiliary; it is not asserted to
be an original equation or to have zero target value.

\subsection{Positive necessity and the exact saving}

Given a solution of the represented system, take its circuit values,
$\delta=1$, all unit copies $X'=X=1$, zero dummy coordinates, and
\[
 u=\left\lceil\frac{x^2}{2}\right\rceil.
\]
The direct guard then has raw value
$2u-x^2+1=1+(x^2\bmod2)$, which is one or two and survives either
incoming carry. All ordinary targets have raw value one, and the
first special target has value one with carry zero. The low
coordinate tests vanish by support separation.

Choose a power-of-two $b$ larger than every coordinate and use the
canonical codes, $B=Hb^2$, $q=B^L$ and $n=q^8$. Their positive
slack witnesses exist as before. In particular, $e_0$ has unit
digit one, the $\delta$ digit makes $g>0$, and the nonconstant
polynomial $\ell_0(T)+e_0(T)T^L$ has nonnegative coefficients, so
its value at $B$ exceeds its value at four and the congruence
quotient $t$ is positive. Also $\Omega=\lambda-e>0$ and
\[
 \mathcal C^2<q\le\theta q,
 \qquad
 S_3=\lambda(\theta q-\mathcal C^2)+e\mathcal C^2>0.
\]
All three masks hold. Their newly calculated $r$ satisfies
\eqref{eq:unit-center-ranges} and the binomial divisibility in
\eqref{eq:unit-center-decoded-powers}.

The positive Pell construction of Section~\ref{sec:composed-99}
now applies with these new coding values. Choose $U=4^{2r+1}$,
$w=U/n^2$, $Y_P=\lfloor(U+1)^{2r}/U^r\rfloor$, and $s=Y_P/n^2$.
Its binomial congruence makes the latter quotient integral.
Choose the main and first Pell coordinates and interval witnesses,
then the second coordinates at $L$, the gap and positive exponent
quotients. In particular
$\Delta=(\psi_{a+4}(L)-\psi_4(L))/a$ is a positive integer.
Finally choose the doubled-index auxiliary witnesses at the new
main parameter. This changes no coding value. Every required
positive unknown is therefore supplied in the same proved order.

The former seven-instruction block computed $2\lambda$, then
$3\lambda=2\lambda+\lambda$. Replace it by the six instructions
\begin{equation}\label{eq:unit-center-six}
 \begin{gathered}
 U_\lambda=\theta\lambda,\qquad
 T_{\mathrm{coef}}=U_\lambda+1,\qquad \lambda_3=3\lambda,\\
 R_2=T_{\mathrm{coef}}+\lambda_3,\qquad
 t_2=\lambda-e,\qquad P_2=U_\lambda q.
 \end{gathered}
\end{equation}
The tests $q^2=R_2$ and $\theta+4=B$ still give exactly
$\lambda(B-1)=q^2-1$. The new block has three multiplications and
three additions; the old block had three multiplications and four
additions. Every other arithmetic instruction is retained.

\begin{theorem}\label{thm:best98}
Every recursively enumerable subset of the positive integers has an
admissible fixed index $(V,H,T_L)$ for a system with 34 positive
unknowns and 22 equations admitting a straight-line certificate of
98 operations: 55 multiplications and 43 additions. Here
$T_L=\psi_4(L)$ for the underlying exponent of the rebuilt index.
The supplied parameters are $(x,V,H,T_L)$, and fixed numerals are free.
\end{theorem}
\begin{proof}
The preceding construction proves the representation in both
directions and gives every required positive witness. Formula
\eqref{eq:unit-center-six} saves one addition from
Theorem~\ref{thm:best99}, giving $55+43=98$.
The exact checker
\texttt{round27\_1980\_unit\_center\_certificate.py} verifies all
98 primitive instructions, all 22 source residuals, and the same
three residual corrections as its predecessor. Only the packed
$r$ equation and the positive-$\sigma$ and positive-$\Omega$
equations change; geometry and every Pell source equation are
unchanged. Its finite regressions are distinct from the universal
proof above. No optimality is asserted.
\end{proof}
